% TeX root = ../main.tex

\chapter{Prelimnaries}
{\color{red} WORK IN PROGRESS!!!}


Thus far, we have been speaking about $C^{\infty}$ differential forms, but we can define differential forms in more general spaces.

\section{Function Spaces of Differential Forms}

Throughout this section, let $\Omega\subset\mathbb R^n$ be an open set and let
$0\leq k\leq n$. Since $\Omega$ is an open subset of $\mathbb R^n$, we identify
$T_p\Omega\simeq\mathbb R^n$ for every $p\in\Omega$. Consequently, every
$differential\ k$-form may be viewed as a map
\[
f:\Omega\longrightarrow\Lambda^k(\mathbb R^n).
\]

With respect to the standard basis of $\Lambda^k(\mathbb R^n)$, every
$k$-form admits a unique expansion of the form
\[
f=\sum_{I\in\mathscr T_k}f^I\,dx^I
=\sum_{1\leq i_1<\cdots<i_k\leq n}
f^{i_1\cdots i_k}\,
dx^{i_1}\wedge\cdots\wedge dx^{i_k},
\]
where each coefficient
$f^I:\Omega\to\mathbb R$ is a real-valued function.

\begin{remark}
Since $\Lambda^k(\mathbb R^n)$ is a finite-dimensional vector space of dimension
$\binom{n}{k}$, a differential $k$-form may equivalently be regarded as a
$\binom{n}{k}$-tuple of scalar-valued functions. All regularity notions are
defined coefficientwise.
\end{remark}

\begin{definition}[$L^p$ differential forms]
Let $1\leq p\leq\infty$. We write
\[
f\in L^p(\Omega;\Lambda^k)
\]
if and only if every coefficient function satisfies
\[
f^I\in L^p(\Omega),\qquad I\in\mathscr T_k.
\]
Equivalently,
\[
L^p(\Omega;\Lambda^k)
=
L^p(\Omega;\mathbb R^{\binom{n}{k}}).
\]
\end{definition}

The pointwise inner product on $\Lambda^k(\mathbb R^n)$ is defined by
\[
\langle f,g\rangle
=
\sum_{I\in\mathscr T_k}f^Ig^I,
\]
with associated pointwise norm
\[
|f|^2=\langle f,f\rangle
=
\sum_{I\in\mathscr T_k}(f^I)^2.
\]

The $L^p$ norm of a differential form is obtained by integrating the pointwise
norm:
\[
\|f\|_{L^p(\Omega;\Lambda^k)}
=
\left(
\int_\Omega |f(x)|^p\,dx
\right)^{1/p},
\]
with the usual modification when $p=\infty$. In particular,
\[
\|f\|_{L^2(\Omega;\Lambda^k)}^2
=
\int_\Omega
\sum_{I\in\mathscr T_k}(f^I)^2\,dx.
\]

\begin{definition}[Sobolev differential forms]
Let $1\leq p<\infty$. We write
\[
f\in W^{1,p}(\Omega;\Lambda^k)
\]
if and only if every coefficient satisfies
\[
f^I\in W^{1,p}(\Omega),
\qquad I\in\mathscr T_k.
\]
Thus a differential form belongs to $W^{1,p}$ precisely when all of its
coefficients possess weak first-order partial derivatives lying in $L^p(\Omega)$.
\end{definition}

For a differential form
$f=\sum_I f^I dx^I$, its gradient is defined coefficientwise by
\[
\nabla f
=
\sum_{i=1}^n\sum_{I\in\mathscr T_k}
\frac{\partial f^I}{\partial x^i}
\,dx^i\otimes dx^I,
\]
and its pointwise norm is
\[
|\nabla f|^2
=
\sum_{i=1}^n
\sum_{I\in\mathscr T_k}
\left(
\frac{\partial f^I}{\partial x^i}
\right)^2.
\]

The Sobolev norm is therefore
\[
\|f\|_{W^{1,p}(\Omega;\Lambda^k)}
=
\left(
\|f\|_{L^p(\Omega;\Lambda^k)}^p
+
\|\nabla f\|_{L^p(\Omega;\Lambda^k)}^p
\right)^{1/p},
\]
and, in particular,
\[
\|f\|_{W^{1,2}(\Omega;\Lambda^k)}^2
=
\|f\|_{L^2(\Omega;\Lambda^k)}^2
+
\|\nabla f\|_{L^2(\Omega;\Lambda^k)}^2.
\]

\begin{definition}[Classical regularity spaces]
For an integer $r\geq0$, we write
\[
f\in C^r(\Omega;\Lambda^k)
\]
if and only if every coefficient function belongs to $C^r(\Omega)$.

Similarly,
\[
f\in C^{r,\alpha}(\Omega;\Lambda^k),
\qquad 0<\alpha\leq1,
\]
if every coefficient belongs to the Hölder space
$C^{r,\alpha}(\Omega)$.
\end{definition}

Recall that the Hölder seminorm of a scalar-valued function $u$ is
\[
[u]_{C^{0,\alpha}(\Omega)}
=
\sup_{\substack{x,y\in\Omega\\x\neq y}}
\frac{|u(x)-u(y)|}{|x-y|^\alpha},
\]
and the $C^{r,\alpha}$ norm is
\[
\|u\|_{C^{r,\alpha}(\Omega)}
=
\sum_{|\beta|\leq r}
\sup_{x\in\Omega}|D^\beta u(x)|
+
\sum_{|\beta|=r}
[D^\beta u]_{C^{0,\alpha}(\Omega)}.
\]
The corresponding norm for differential forms is again defined
componentwise.

\begin{definition}[Spaces on the closure]
We write
\[
f\in C^r(\overline{\Omega};\Lambda^k)
\]
if every coefficient extends continuously to $\overline{\Omega}$ together
with all derivatives of order at most $r$.

Likewise,
\[
f\in C^{r,\alpha}(\overline{\Omega};\Lambda^k)
\]
if every coefficient extends to a function in
$C^{r,\alpha}(\overline{\Omega})$.
\end{definition}

\begin{remark}
The distinction between $C^{r,\alpha}(\Omega)$ and
$C^{r,\alpha}(\overline{\Omega})$ is important. Functions in
$C^{r,\alpha}(\Omega)$ are only required to satisfy the Hölder condition in the
interior of $\Omega$, whereas elements of
$C^{r,\alpha}(\overline{\Omega})$ possess continuous extensions, together with
their derivatives up to order $r$, to the boundary $\partial\Omega$. Boundary
conditions such as $\nu\wedge f=0$ or $\nu\lrcorner f=0$ are naturally stated in
the latter setting.
\end{remark}

\begin{remark}
Since $\Lambda^k(\mathbb R^n)$ is finite dimensional, all norms on
$\Lambda^k(\mathbb R^n)$ are equivalent. Consequently, the particular choice of
Euclidean norm used to define the spaces above is immaterial, and all resulting
$L^p$, Sobolev and Hölder norms are equivalent.
\end{remark}

%this is where serious business is there i suppose
%i think
